\appendix
\section{Appendix A: Mathematical Foundations of the Game Formalism}
\label{appendix:math_foundations}

This appendix provides formal mathematical background and complete reference derivations behind the specifications introduced in the main text.

\subsection{Run Semantics as an MDP over Intervention Spaces}
As described in Sec.~\ref{sec:formaldefinition}, each game instance is a causal MDP:
\[
    \mathcal{M}
    = \bigl(\statespace, \actionspace, \transition, \phi, \resultmetric, \behaviormetric \bigr),
\]
where $\statespace$ is a \textbf{history-valued} state space
and each action $a \in \actionspace$ is either
(i) a batch of interventions $(\intervention, n)$ or
(ii) a final answer action.
We emphasize that the environment is \emph{fully deterministic given the SCM and seed} — the only randomness comes from SCM sampling.

The run $r \in \runs$ is thus defined as:
\[
    r = (s_0, a_0, s_1, a_1, \dots, s_T),
    \qquad s_{t+1} = \transition(s_t, a_t),
\]
where $T$ is either agent-chosen (final action) or forced (via $\phi$).
The transcript is the persistent serialization of $(r, \text{answer})$.

\subsection{Pareto Interpretation of Dual Metrics}
Unlike reward-maximizing MDPs, \game deliberately does not collapse to a single scalar reward.  
Instead, final evaluation yields:
\[
    \Bigl( \resultmetric(\text{answer}),\ \behaviormetric(r) \Bigr)
    \in \mathbb{R}^2,
\]
reflecting the \textbf{two fundamental desiderata} of causal cognition:
\begin{itemize}
    \item correctness of inference (judged externally via ground truth),
    \item resource-efficiency or experimental prudence.
\end{itemize}
Whether or not these are later combined (e.g., for agent ranking) is left to the researcher.

\bigskip
Additional technical derivations, SCM generation proofs, hooks API, and DTO schemas may be found in the subsequent appendices.


\section{Appendix B: DTO Schemas and Information-Theoretic Guarantees}
\label{appendix:dtos}

A key architectural principle of \game is the strict use of \textbf{Data Transfer Objects (DTOs)} — immutable, minimal, transportation-safe containers of game state.
They enforce \textbf{type safety}, \textbf{role restriction}, and avoid accidental leakage of internal SCM information to the agent.

\subsection{Core DTOs Exchanged During Runtime}

\paragraph{TranscriptEntryDTO.}
Represents a single round of interaction — contains only:
\begin{itemize}
    \item the intervention decision requested by the agent (as a typed, validated object),
    \item the interventional samples returned by the environment (never more than allowed),
    \item optional feedback or metadata (e.g.\ partial evaluation score, if mission allows),
    \item a budget snapshot from the BudgetManager.
\end{itemize}
It is \textbf{impossible} for this DTO to accidentally contain or leak SCM internals such as $f_i$, $\eta_i$, or the true DAG.

\paragraph{SamplesDTO.}
Encodes interventional data as a \textbf{mapping from intervention $\intervention$ to sampled dataset}:
\[
    \texttt{SamplesDTO} \;:\;\;
    \intervention \mapsto \{(x^{(j)}, y^{(j)})_{j=1}^n\}.
\]
All samples are drawn through the SCM equation engine under seed control.  
No symbolic form of the SCM is ever exposed — only \textbf{realized outcomes}.

\paragraph{AnswerDTO.}
Returned only when the agent stops — contains:
\begin{itemize}
    \item predicted DAG $\,\hat{\DAG}$, or
    \item estimated treatment-effect function, or
    \item reconstructed SCM parameters — depending on the mission type.
\end{itemize}
This is the only node of communication explicitly judged by $\resultmetric$.

\subsection{Security and Information Containment Guarantees}

\begin{itemize}
    \item DTOs enforce a \textbf{one-way visibility barrier} — the agent \emph{only sees data consistent with the contractual reality of the game}.
    \item The agent \emph{cannot observe} latent variables, hidden SCM functions $f_i$, or the full universe of possible interventions, unless defined so in the Spec.
    \item Every DTO crossing from environment to agent is automatically vetted for \textbf{leak safety}.
\end{itemize}

This ensures rigorous \textbf{causal identifiability testing} — agents must legitimately \textbf{infer causality from experimental evidence}, never from accidental metadata exposure.

\bigskip
Following this, Appendix C will detail the \textbf{Hook API and Budget Enforcement Mechanism}, which control the flow of time, stopping logic, and lifecycle events in \game.


\section{Appendix C: Hooks, Lifecycle, and Budget Enforcement}
\label{appendix:hooks_budget}

In this section, we provide a precise specification of the event lifecycle within \game, emphasizing the role of \textbf{HookManager} and \textbf{BudgetManager} in controlling execution, termination, and observability.

\subsection{Event Lifecycle and Hook Triggers}
Hooks provide \emph{instrumentation without interference} — they can observe events but cannot alter decisions or leak SCM internals to the agent.

\paragraph{Canonical Hook Events.}
\begin{itemize}
    \item \textbf{GAME\_START} — triggered before the first round.
    \item \textbf{ROUND\_START} — triggered before agent observes the updated state.
    \item \textbf{ROUND\_END} — after samples are returned but before stop check.
    \item \textbf{GAME\_END} — after the agent stops or BudgetManager forces termination.
\end{itemize}

These hooks enable logging, plotting, telemetry, debugging, and human visualization — yet \textbf{never override agent or SCM decisions}.

\subsection{Budget Enforcement}
BudgetManager tracks consumption over time. Let $\beta_t$ represent current budget usage at round $t$.

\begin{itemize}
    \item Costs may be assigned per round, per intervention, per sample, or any combination.
    \item A strict constraint $\beta_t \leq \beta_{\max}$ must hold at all times.
    \item If $\beta_t > \beta_{\max}$, the game is \emph{immediately stopped} via forced termination.
\end{itemize}

\subsection{Forced vs.\ Voluntary Termination}
The game ends in exactly one of two ways:
\begin{itemize}
    \item \textbf{Voluntary Stop} — the agent explicitly issues a valid stop action and returns an AnswerDTO.
    \item \textbf{Forced Stop} — BudgetManager detects budget exhaustion before the agent stops.
\end{itemize}

In the forced case, no final answer is collected — $\resultmetric$ is applied to an “empty answer,” ensuring the agent is \textbf{penalized for overshooting resources}.

\bigskip
Having described DTO boundaries, lifecycle events, and budget rules, we now finalize the architecture with the exact \textbf{Spec Contracts and Registry Mechanism} used to dynamically instantiate all components of the game.


\section{Appendix D: Specification Contracts and Dynamic Registry}
\label{appendix:specs_registry}

All components of \game — Agents, SCMs, Missions, Metrics, Hooks, etc. — are instantiated \textbf{declaratively} via \textbf{Spec objects}.

\subsection{Spec Contract}
Each Spec is a typed, schema-validated configuration object containing:
\begin{itemize}
    \item \texttt{class\_} — a string import path (e.g. \texttt{"cg.agents.RandomAgent"})
    \item \texttt{id} — an informational label
    \item \texttt{params} — a JSON-serializable dictionary of arguments
\end{itemize}

The Spec itself does not execute — it is \textbf{purely descriptive}.  
This enables perfect static reproducibility independent of runtime environment.

\subsection{Registry: From Spec to Runtime Instance}
Upon calling \texttt{build(spec)}, the Registry:
\begin{enumerate}
    \item dynamically imports the Python class specified by \texttt{class\_},
    \item validates the \texttt{params} against the Contract interface,
    \item returns a live instance implementing the exact abstract \textbf{Contract}
          (e.g.\ \texttt{Agent}, \texttt{SCM}, \texttt{Metric}).
\end{enumerate}

The Registry thus enforces:
\begin{itemize}
    \item \textbf{decoupling} between configuration and execution logic,
    \item \textbf{pluggability} — arbitrary user code may be injected if it fulfills the interface,
    \item \textbf{complete reproducibility} — behavior fully determined by Spec + seed.
\end{itemize}

\subsection{Safe Extensibility for Benchmark Standardization}
Crucially, multiple researchers can share \textbf{identical Spec files} but supply different agent implementations —  
or contribute official Specs to converge community consensus on:
\begin{itemize}
    \item standard benchmark missions,
    \item official SCMs,
    \item agreed-upon budget regimes.
\end{itemize}

This Spec–Contract–Registry system is the foundational mechanism that enables \textbf{benchmark standardization without code centralization}.

\bigskip
In Appendix~E we finalize with **reproducibility guarantees**, handling of random seeds, and practical execution guidelines.


\section{Appendix E: Reproducibility, Seeding, and Execution Guarantees}
\label{appendix:reproducibility}

A central design principle of \game is that \textbf{every experiment is perfectly reproducible}, independent of machine, hardware, OS, or execution order — as long as the same \textbf{Spec + seed} are used.

\subsection{Global Reproducibility Contract}
Every official \game run is defined uniquely by:
\[
    \bigl(\texttt{Spec},\ \texttt{seed} \bigr).
\]
This includes the SCM choice, the agent configuration, mission metrics, start/stop rules, and hook definitions.

The seed determines all stochasticity in:
\begin{itemize}
    \item SCM sampling (e.g.\ noise terms, categorical resolution),
    \item environment transitions,
    \item agent random exploration (if applicable).
\end{itemize}

\subsection{Seed Handling Strategy}
\begin{itemize}
    \item The \textbf{global seed} is fixed at game construction time.
    \item All internal RNGs — SCM sampling, interventional resampling, agent decisions — derive from this seed in a controlled hierarchy.
    \item Executing the same Spec + seed on different machines must produce bitwise identical Transcripts and evaluation scores.
\end{itemize}

\subsection{Determinism Across Parallel Execution}
Even if many game instances are executed in parallel via the Runner, each one receives its own isolated seed context.  
\textbf{Parallel execution does not alter execution order or outcomes} — it only accelerates the batch.

\subsection{Official Benchmark Integrity}
This guarantees that:
\begin{itemize}
    \item any officially shared result can be \textbf{exactly reproduced} by any other party,
    \item no difference in hardware or multiprocessing backend can compromise fairness,
    \item benchmark releases in \game may be treated as \textbf{permanent reference artifacts},
    \item contributions from multiple institutions can be compared \textbf{without trust assumptions}.
\end{itemize}

\bigskip
This concludes the reproducibility and execution guarantees.  
Appendix~F will provide deeper details on configuration, resource management, and advanced customization for expert users.

\section{Appendix F: Advanced Configuration and Customization}
\label{appendix:advanced}

This appendix is intended for expert users who wish to push \game beyond its default configuration,
including advanced SCM construction, agent-side internal memory control, and highly customized mission definitions.

\subsection{Fine-Grained SCM Configuration}
In all three SCM generation pathways (dataset-based, Bayesian-network–derived, or fully synthetic), 
\game supports advanced constraints — for example:
\begin{itemize}
    \item \textbf{topological depth control} (min/max causal path length),
    \item \textbf{heterogeneous variable typing} (numerical vs.\ categorical),
    \item \textbf{nonlinear functional syntax grammars with user-imposed limitations} (polynomial degree, trigonometric allowance),
    \item \textbf{custom intervention availability schedules} (e.g., only certain variables become controllable after $t{=}5$).
\end{itemize}
Such parameters may be declared directly inside the SCMSpec, making SCM family inheritance and fine-grained experimental variation trivial.

\subsection{Agent Memory and Internal State}
While the Transcript only exposes limited observable information, \emph{agents are free to internally maintain arbitrary memory}, including:
\begin{itemize}
    \item Bayesian belief states,
    \item learned causal graph hypotheses,
    \item experimental design policies that adapt as evidence accumulates.
\end{itemize}
The only restriction is that this memory must be private — it may not influence the game’s external machinery (i.e., cannot bypass DTO boundaries).

\subsection{Mission-Specific Customization}
Beyond the standard missions (graph discovery, CATE estimation, SCM reconstruction), researchers may define:
\begin{itemize}
    \item novel \resultmetric{} with arbitrary semantic meaning,
    \item multi-objective or lexicographically ordered dominance rules,
    \item domain-restricted forms of counterfactual queries (e.g., multi-variable nested do-operations).
\end{itemize}
These extensions naturally integrate into the game, as long as the \textbf{Mission Contract} is implemented.

\subsection{Controlled Environment Modifications}
\game supports controlled configuration of environment behavior, such as:
\begin{itemize}
    \item \textbf{forced burn-in samples} before agent control begins,
    \item \textbf{dynamic action availability} (actions appear/disappear across time),
    \item \textbf{curriculum-style progressive complexity increase} across batched Runner calls.
\end{itemize}

\bigskip
All such customization is possible \emph{without modifying the \game core}, validating its modularity and extensibility for research-grade experimentation.
Researchers are encouraged to register official open benchmark variants via shared Spec repositories.

Appending further, Appendix~G will contain fully worked example configurations and code-level usage templates.


\section{Appendix G: Complete Example Configurations (Specs)}
\label{appendix:example_specs}

To make the abstract architecture concrete, this appendix provides example \textbf{fully executable Spec configurations}
for different mission types — ready to be loaded directly via the Registry.

Each Spec is a \textbf{declarative JSON-like structure} (pseudocode below, compatible with Pydantic or \texttt{BaseModel}-style construction).  
These examples illustrate minimal but valid configurations that obey all Contract rules.

\subsection{Example 1 — DAG Discovery Mission on a Physical Law SCM}
```json
{
  "game_id": "ideal_gas_dag_inference",
  "scm": {
    "class_": "cg.scms.PhysicalLawSCM",
    "id": "ideal_gas",
    "params": { "law": "ideal_gas" }
  },
  "mission": {
    "class_": "cg.missions.DAGInferenceMission",
    "id": "dag_recovery",
    "params": { "metric": "SHD" }
  },
  "stopping": {
    "class_": "cg.stopping.RoundLimit",
    "id": "round_cap_20",
    "params": { "max_rounds": 20 }
  },
  "agents": [
    { "class_": "cg.agents.RandomAgent", "id": "random_1", "params": { "seed": 101 } }
  ]
}
```

\section{Appendix H: Transcript Inspection and Post-Processing}
\label{appendix:transcript_postprocessing}

This appendix demonstrates how to \textbf{programmatically inspect, analyze, and extend} the outputs of a completed \game run — in particular the Transcript, which acts as the immutable record of all agent–environment interaction.

\subsection{Loading a Run}
After executing:
\begin{minted}[fontsize=\small]{python}
runs = game.run()
\end{minted}
the variable \texttt{runs} contains one Transcript per agent.  
Then, for example:
\begin{minted}[fontsize=\small]{python}
r = runs[0]  # a single agent run
print(r.summary())
\end{minted}

Typical fields include:
\begin{itemize}
    \item number of rounds executed,
    \item total samples collected,
    \item stop reason (voluntary vs.\ budget-forced),
    \item final result (AnswerDTO) if agent stopped voluntarily.
\end{itemize}

\subsection{Round-by-Round Inspection}
You may inspect the Transcript step-by-step:
\begin{minted}[fontsize=\small]{python}
for entry in r.entries:
    print("Intervention:", entry.decision)
    print("Samples returned:", entry.samples)
    print("Budget snapshot:", entry.budget_snapshot)
\end{minted}

This is fully deterministic — guaranteed identical across machines given the same Spec + seed.

\subsection{Custom Metric Post-Processing}
Beyond the official \resultmetric{}, analysts can run custom evaluations:
\begin{minted}[fontsize=\small]{python}
custom_score = custom_evaluator(r)
``\end{minted}

Since the Transcript is immutable, this never affects the official evaluation, preserving benchmark fairness.

\subsection{Optional Export}
Transcripts can be safely exported to JSON for third-party plotting or statistical analysis:
\begin{minted}[fontsize=\small]{python}
r.to_json("experiment_run_01.json")
\end{minted}

\bigskip
Transcript inspection is foundational to auditability, fairness, and interpretability.  
In Appendix~I, we present demonstrative end-to-end workflows combining \game with external scientific tools (e.g.\ PyTorch, symbolic solvers, or LLM copilots).

\section{Appendix J: Practical Implementation Notes and Contributor Guidance}
\label{appendix:implementation_notes}

This appendix summarizes pragmatic advice for researchers and contributors wishing to extend or maintain \game effectively.

\subsection{Recommended Project Structure}
When extending with new agents, missions, or SCMs, we recommend preserving the standard modular folder structure:
\begin{itemize}
    \item \texttt{contracts/} — defines core interfaces
    \item \texttt{specs/} — declarative configurations only
    \item \texttt{dto/} — strictly serialized data objects
    \item \texttt{runtime/} — environment + orchestration logic
    \item \texttt{agents/} — agent implementations
    \item \texttt{metrics/} — behavior / result metric classes
    \item \texttt{scms/} — SCM generators or wrappers
    \item \texttt{hooks/} — optional logging / analysis utilities
\end{itemize}

\subsection{Contribution Best Practices}
To preserve reproducibility and safety:
\begin{itemize}
    \item Always implement against the appropriate \textbf{Contract} (e.g.\ \texttt{Agent}, \texttt{SCM}, \texttt{Mission}).
    \item Never bypass DTOs to expose internal state — all agent inputs must be mediated through the runtime.
    \item Validate parameters at Spec instantiation time wherever possible.
    \item Include minimal runnable examples for every new component.
\end{itemize}

\subsection{Testing and Validation}
We recommend:
\begin{itemize}
    \item unit tests for contract compliance,
    \item snapshot tests to confirm reproducibility under fixed seeds,
    \item short parallelized runs to ensure no concurrency-induced non-determinism,
    \item optional integration tests for cross-module compatibility.
\end{itemize}

\subsection{Migration Stability}
The API is designed to remain stable — even as missions or agents evolve.  
When modifying core contracts, ensure all changes are backward-compatible with existing Specs.
If breaking changes are unavoidable:
\begin{itemize}
    \item Introduce versioned contracts (e.g., \texttt{AgentV2}),
    \item Provide automated migration tools or adapters.
\end{itemize}

\bigskip
This marks the end of the appendices.  
The \game framework is built to sustain long-term, community-driven evolution of causal inference evaluation — rigorously, transparently, and reproducibly.


\section{Appendix K: Table of Symbols and Notation}
\label{appendix:table_of_symbols}

\begin{table}[h!]
\centering
\renewcommand{\arraystretch}{1.15}
\begin{tabular}{ll}
\toprule
\textbf{Symbol} & \textbf{Meaning} \\
\midrule
$\scm$ & Underlying Structural Causal Model \\
$\varset$ & Set of all variables in the SCM \\
$\latentx$ & Set of latent (non-observable) variables \\
$\measuredx$ & Set of observed (but not controllable) variables \\
$\controllablex$ & Set of controllable (intervenable) variables \\
$\interventions$ & Space of all admissible intervention assignments \\
$\intervention$ & A specific intervention $\intervention : \varname \mapsto \text{value} \cup \{\none\}$ \\
$\experiments$ & Set of batched intervention requests $e : \interventions \to \mathbb{N}$ \\
$\dataset$ & Dataset associated with a particular intervention, i.e.\ observations \\
$\statespace$ & Set of all environment states \\
$\actionspace$ & All admissible agent actions (interventions or stop decision) \\
$\transition$ & Environment transition function $\transition : \statespace \times \actionspace \to \statespace$ \\
$\initstate$ & Initial state at the beginning of the game \\
$\runs$ & Set of all possible state--action trajectories \\
$\terminalstate$ & Terminal absorbing state after voluntary or forced stop \\
$\goal$ & Mission statement / objective specification \\
$\outputspace$ & Space of answer objects an agent may return at stop \\
$\resultmetric$ & External metric evaluating correctness of the returned answer \\
$\behaviormetric$ & External metric evaluating cost-efficiency / experimental behavior \\
$\DAG$ & Directed acyclic graph induced by the SCM \\
$\edgesset$ & Set of directed causal edges in the DAG \\
$\none$ & Special symbol indicating no intervention on a variable \\
\bottomrule
\end{tabular}
\caption{Canonical notation used throughout the \game formalism.}
\end{table}
